%% kapa: exam=1 round=25 year=2014 period=2 form=A subject=appraisal_law subject_ko=부동산관계법규 seq=30 number=30
%% source: Q-Net 감정평가사 기출문제 (cst00309) artlSeq=5101207 file=extracted/1008/5101207/2014년 제25회 감정평가사 1차시험 기출문제/2교시(A형).hwp
%% conversion: kordoc 4.24.51 -> tools/gen1.py (auto); review=proofread-sonnet; flags=none
\begin{question}{30}
건축법령상 건축 관련 입지와 규모의 사전결정에 관한 설명으로 옳지 않은 것은?
\begin{choices}
  \item 건축허가 대상 건축물을 건축하려는 자는 건축허가를 신청하기 전에 허가권자에게 그 건축물을 해당 대지에 건축하는 것이「건축법」이나 다른 법령에서 허용되는지에 대한 사전결정을 신청할 수 있다.
  \item 사전결정신청자는 건축위원회 심의와「도시교통정비 촉진법」에 따른 교통영향분석ㆍ개선대책의 검토를 동시에 신청할 수 있다.
  \item 허가권자는 사전결정이 신청된 건축물의 대지면적이「환경영향평가법」에 따른 소규모 환경영향평가 대상사업인 경우 환경부장관이나 지방환경관서의 장과 소규모 환경영향평가에 관한 협의를 하여야 한다.
  \item 사전결정 통지를 받은 경우에도「국토의 계획 및 이용에 관한 법률」에 따른 개발행위허가는 따로 받아야 한다.
  \item 사전결정신청자가 사전결정을 통지받은 날부터 2년 이내에 건축허가를 신청하지 아니하면 사전결정의 효력이 상실된다.
\end{choices}
\end{question}
\answer{4}
